
Supposons en particulier que \(H=\Ghat\)~; alors le \(n\)-ième polynôme universel de \(\Phi\) n'a pas de composantes de poids \(<n\). Considérons maintenant le morphisme canonique~: \(H\longrightarrow1+\widehat{\operatorname{Gr}\,\mathord{.}\,}^{+}\)~; par composition avec \(\Phi\), on obtient un morphisme~: \((\Ghat)^k\longrightarrow1+\widehat{\operatorname{Gr}\,\mathord{.}\,}^{+}\), qui d'après la démonstration de 1.12 se factorise par \((1+\widehat{\operatorname{Gr}\,\mathord{.}\,}^{+})^k\) et définit un morphisme \(\Phi:(1+\widehat{\,\mathord{.}\,}^{+})^k\longrightarrow1+\widehat{\,\mathord{.}\,}^{+}\)~; on voit aussitôt que le \(n\)-ième polynôme universel de \(\Phi\) est le même que celui de \((\Ghat)^k\longrightarrow1+\widehat{\operatorname{Gr}\,\mathord{.}\,}^{+}\), c'est-à-dire la composante isobare de poids \(n\) du \(n\)-ième polynôme universel de \(\Phi\).

Les mêmes remarques s'appliquent au cas \(H=\Ga\)~: on peut définir un morphisme de foncteurs \(H\longrightarrow\widehat{\operatorname{Gr}\,\mathord{.}\,}^{+}\) (avec \(\widehat{\operatorname{Gr}\,\mathord{.}\,}^{+}=\widehat{\,\mathord{.}\,}^{+}\circ\operatorname{Gr}\), \(\widehat{\,\mathord{.}\,}^{+}\) étant le foncteur défini sur \(\mathcal G\) par \(\widehat B^+=\prod_{i\geq1}B^i\)) de façon analogue à 1.11; on recopie alors raisonnement et conclusions de l'alinéa précédent.

1.16. Terminons sur un exemple. On suppose que \(K=Q\). On définit un endomorphisme du foncteur \(1+\widehat{\,\mathord{.}\,}^{+}\), noté \(\Toddc\), et appelé caractère de Todd, en posant~:
\[
  \Toddc(1+X)=\frac{X}{1-e^{-X}}.
\]
Les conditions i) et ii) du théorème 1.13 se vérifient immédiatement, ce qui montre que nous avons bien défini ainsi un endomorphisme de \(1+\widehat{\,\mathord{.}\,}^{+}\).

\section*{2. Définition des \(\lambda\)-anneaux; exemples}

Dans ce paragraphe, l'anneau de base du paragraphe 1 sera \(\mathbb Z\) (sauf en 2.9)~; on se place sur la catégorie \(\mathcal C_0\) (1.8) et on note \(\Ghat\) la restriction de \(\Ghat\) à \(\mathcal C_0\) (1.10).

